\documentclass[12pt, a4paper]{report}

% Packages :

\usepackage[french]{babel}
\usepackage[utf8]{inputenc}
\usepackage[T1]{fontenc}
\usepackage[pdftex, pdfauthor={Bacomathiques}]{hyperref}
\usepackage{sectsty}
\usepackage[explicit]{titlesec}
\usepackage{xcolor}
\usepackage{amsmath}
\usepackage{amssymb}
\usepackage{amsthm}
\usepackage{fourier}
\usepackage{titlesec}
\usepackage{fancyhdr}
\usepackage{catchfilebetweentags}
\usepackage[french, capitalise, noabbrev]{cleveref}
\usepackage[fit, breakall]{truncate}
\usepackage[margin=3cm]{geometry}
\usepackage{tocloft}
\usepackage{tikz}
\usepackage{tocloft}
\usepackage{microtype}
\usepackage{listings}
\usepackage{tabularx}
\usepackage{calc}
\usepackage[export]{adjustbox}
\usepackage[most]{tcolorbox}
\usepackage{standalone}
\usepackage{xlop}
\usepackage{etoolbox}
\usepackage{environ}

\usetikzlibrary{arrows.meta}
\usetikzlibrary{trees}

% Paramètres :

\author{Bacomathiques}
\definecolor{graphe}{HTML}{93c9ff}
\setcounter{MaxMatrixCols}{12}
\setlength{\parindent}{0pt}
\setlength{\fboxsep}{0pt}
%\pdfsuppresswarningpagegroup=1

% Code :

\lstdefinestyle{style}{
	backgroundcolor=\color{white},
	commentstyle=\em\color[HTML]{999988},
	keywordstyle=\bfseries,
	identifierstyle=\normalfont,
	stringstyle=\color[rgb]{0.87, 0.07, 0.27},
	basicstyle=\ttfamily\color{black},
	breakatwhitespace=false,
	breaklines=true,
	captionpos=b,
	keepspaces=true,
	numbers=left,
	numbersep=5pt,
	showspaces=false,
	showstringspaces=false,
	showtabs=false,
	tabsize=2,
	numbers=none
}

\lstset{style=style}
\lstset{
	literate=
	{á}{{\'a}}1
	{à}{{\`a}}1
	{ã}{{\~a}}1
	{é}{{\'e}}1
	{ê}{{\^e}}1
	{í}{{\'i}}1
	{ó}{{\'o}}1
	{õ}{{\~o}}1
	{ú}{{\'u}}1
	{ü}{{\"u}}1
	{ç}{{\c{c}}}1
}

\lstset{
	framextopmargin=10pt,
	framexrightmargin=10pt,
	framexbottommargin=10pt,
	framexleftmargin=10pt,
	xleftmargin=10pt,
	xrightmargin=10pt,
}

% Couleurs :

\definecolor{title}{HTML}{912c21}
\definecolor{section}{HTML}{1c567d}
\definecolor{subsection}{HTML}{2980b9}

\definecolor{rule}{HTML}{c4c4c4}

\definecolor{formula}{HTML}{ebf3fb}
\definecolor{formula-left}{HTML}{3583d6}

\definecolor{tip}{HTML}{dcf3d8}
\definecolor{tip-left}{HTML}{26a65b}

\definecolor{demonstration}{HTML}{fff8de}
\definecolor{demonstration-left}{HTML}{f1c40f}

\definecolor{exercise}{HTML}{e0f2f1}
\definecolor{exercise-left}{HTML}{009688}

\definecolor{correction}{HTML}{e0f7fa}
\definecolor{correction-left}{HTML}{00bcd4}

\definecolor{toc}{HTML}{fceae9}
\definecolor{toc-left}{HTML}{e74c3c}
\definecolor{toc-dark}{HTML}{87281f}

% Titres :

\renewcommand{\thesection}{\Roman{section} - }
\renewcommand{\thesubsection}{\arabic{subsection}. }

\newcommand{\sectionstyle}{\normalfont\LARGE\bfseries\color{section}}
\titleformat{\section}{\sectionstyle}{\thesection #1}{0pt}{}
\titleformat{name=\section, numberless}{\sectionstyle}{#1}{0pt}{}

\newcommand{\subsectionstyle}{\normalfont\Large\bfseries\color{subsection}}
\titleformat{\subsection}{\subsectionstyle}{\thesubsection #1}{0pt}{}
\titleformat{name=\subsection, numberless}{\subsectionstyle}{#1}{0pt}{}

\titlelabel{\thetitle\ }

% Table des matières :

\addto\captionsfrench{\renewcommand\contentsname{}}
\renewcommand{\cftsecpagefont}{\color{toc-dark}}
\renewcommand{\cftsubsecpagefont}{\color{toc-dark}}
\renewcommand{\cftsecleader}{\cftdotfill{\cftdotsep}}
\renewcommand{\cftsecfont}{\bfseries}
\renewcommand{\cftsecpagefont}{\bfseries\color{toc-dark}}
\setlength{\cftbeforetoctitleskip}{0pt}
\setlength{\cftaftertoctitleskip}{0pt}
\setlength{\cftsecindent}{0pt}
\setlength{\cftsubsecindent}{20pt}
\setlength{\cftsubsecnumwidth}{20pt}

% Commandes :

\newcommand{\newpar}{\\[\medskipamount]}
\newcommand{\lesson}[3]{%
	\newcommand{\level}{#1}%
	\newcommand{\id}{#2}%
	\hypersetup{pdftitle={#3}}
	\begin{center}%
		\includegraphics[width=150px]{\imagespath/bacomathiques}%
		
		\vspace{30pt}%
		{\Huge\color{title} #3}%
		
		\vspace{10pt}%
		{Bacomathiques --- \href{https://bacomathiqu.es/cours/#1/#2}{\color{section} https://bacomathiqu.es}}%
		
		\vspace{20pt}%
	\end{center}%
	\begin{toc}
		\tableofcontents%
	\end{toc}
	\thispagestyle{empty}%
	\newpage%
	\setcounter{page}{1}%
}
\newcommand{\imagespath}{../../images}
\newcommand{\lessonimagespath}{\imagespath/lessons/\level/\id/}
\newcommand{\includelatexpicture}[2][\textwidth - 100pt]{%
	\begin{center}%
		\resizebox{#1}{!}{%
			\input{\lessonimagespath#2}%
		}%
	\end{center}%
	\medskip%
}
\newcommand{\includeimage}[1]{%
	\begin{center}%
		\includegraphics{\lessonimagespath#1}%
	\end{center}%
	\medskip%
}
\newcommand{\includerepresentation}[1]{%
	\begin{center}%
		\setlength{\fboxrule}{0.5pt}%
		\href{https://www.geogebra.org/m/#1}{\includegraphics[width=\textwidth-1pt,fbox]{\lessonimagespath#1}}%
	\end{center}%
}
\newcommand{\floor}[1]{\lfloor #1 \rfloor}

\makeatletter
\newcommand\inputcontent{\@ifstar{\inputcontent@star}{\inputcontent@nostar}}
\newcommand{\inputcontent@star}[1]{%
	\ExecuteMetaData[#1]{content}%
}
\newcommand{\inputcontent@nostar}[1]{%
	\newpage%
	\inputcontent@star{#1}%
}
\makeatother

\let\oldsection\section
\renewcommand\section{\clearpage\oldsection}
\newcommand{\contentwidth}[1][medium]{}

% En-têtes :

\pagestyle{fancy}

\renewcommand{\sectionmark}[1]{\markboth{\thesection \ #1}{}}

\fancyhead[R]{\truncate{0.23\textwidth}{\color{title}\thepage}}
\fancyhead[L]{\truncate{0.73\textwidth}{\color{title}\leftmark}}
\fancyfoot[C]{\scriptsize \href{https://bacomathiqu.es/cours/\level/\id}{\texttt{bacomathiqu.es}}}

\makeatletter
\patchcmd{\f@nch@head}{\rlap}{\color{rule}\rlap}{}{}
\patchcmd{\headrule}{\hrule}{\color{rule}\hrule}{}{}
\makeatother

% Environnements :

\newenvironment{nosummary}{}{}
\newcommand{\tcolorboxtitle}[2]{\setlength{\fboxsep}{2.5pt}\hspace{-10pt}\colorbox{#1-left}{\hspace{8pt}\MakeUppercase{#2} \hspace{2pt} \includegraphics[height=0.8em]{\imagespath/bubbles/#1}\hspace{5pt}}}
\newcommand{\tcolorboxsubtitle}[2]{\ifstrempty{#2}{}{\textcolor{#1-left}{\large#2}\\[\medskipamount]}}
\tcbset{
	frame hidden,
	boxrule=0pt,
	boxsep=0pt,
	enlarge bottom by=8.5pt,
	enhanced jigsaw,
	boxed title style={sharp corners,boxrule=0pt,coltitle={white},titlerule=0pt},
	fonttitle=\fontsize{6pt}{6pt}\bfseries\boldmath,
	top=10pt,
	right=10pt,
	bottom=10pt,
	left=10pt,
	arc=0pt,
	outer arc=0pt,
}
\newtcolorbox{toc}[1][]{
	colback=toc,
	borderline west={3pt}{0pt}{toc-left},
	title=\tcolorboxtitle{toc}{Table des matières},
	colbacktitle=toc,
	before upper={\tcolorboxsubtitle{toc}{#1}}
}
\newtcolorbox{formula}[1][]{
	colback=formula,
	borderline west={3pt}{0pt}{formula-left},
	title=\tcolorboxtitle{formula}{À retenir},
	colbacktitle=formula,
	before upper={\tcolorboxsubtitle{formula}{#1}}
}
\newtcolorbox{tip}[1][]{
	colback=tip,
	borderline west={3pt}{0pt}{tip-left},
	title=\tcolorboxtitle{tip}{À lire},
	colbacktitle=tip,
	before upper={\tcolorboxsubtitle{tip}{#1}}
}
\newtcolorbox{demonstration}[1][]{
	colback=demonstration,
	borderline west={3pt}{0pt}{demonstration-left},
	title=\tcolorboxtitle{demonstration}{Démonstration},
	colbacktitle=demonstration,
	before upper={\tcolorboxsubtitle{demonstration}{#1}}
}

\NewEnviron{whitetabularx}[1]{%
	\renewcommand{\arraystretch}{2.5}
	\colorbox{white}{%
		\begin{tabularx}{\textwidth}{#1}%
			\BODY%
		\end{tabularx}%
	}%
}

% Longueurs :

\newlength{\espacetitreliste}
\setlength{\espacetitreliste}{-16pt}
\newcommand{\entretitreetliste}{\vspace{\espacetitreliste}}

\begin{document}
	%<*content>
	\lesson{premiere}{geometrie-reperee}{Chapitre VI -- Géométrie repérée}

	\section{Le produit scalaire}

	\subsection{Définition}
	\label{definition}

	\begin{formula}[Définition]
		Soient $\overrightarrow{u} = \begin{pmatrix} {x_1} \\ {y_1}\end{pmatrix}$ et $\overrightarrow{v} = \begin{pmatrix} {x_2} \\ {y_2}\end{pmatrix}$ deux vecteurs du plan (c'est-à-dire possédant chacun deux coordonnées).
		\newpar
		Le \textbf{produit scalaire} entre $u$ et $v$, noté $\overrightarrow{u} \cdot \overrightarrow{v}$ est le réel suivant :
		\newpar
		$\overrightarrow{u} \cdot \overrightarrow{v} = x_1 x_2 + y_1 y_2$.
	\end{formula}

	\begin{formula}[Propriétés]
		Soient $\overrightarrow{u}$, $\overrightarrow{v}$ et $\overrightarrow{w}$ des vecteurs du plan et $\lambda \in \mathbb{R}$, on a les propriétés suivantes :
		\begin{itemize}
			\item $\overrightarrow{u} \cdot \overrightarrow{v} = \overrightarrow{v} \cdot \overrightarrow{u}$
			\item $\overrightarrow{u} \cdot (\overrightarrow{v} + \overrightarrow{w}) = \overrightarrow{u} \cdot \overrightarrow{v} + \overrightarrow{u} \cdot \overrightarrow{w}$
			\item $\lambda(\overrightarrow{u} \cdot \overrightarrow{v}) = (\lambda \overrightarrow{u}) \cdot \overrightarrow{v} = \overrightarrow{u} \cdot (\lambda \overrightarrow{v})$
		\end{itemize}
	\end{formula}

	À l'aide du produit scalaire, il est possible de calculer la \textbf{norme} d'un vecteur.

	\begin{formula}[Calcul de la norme]
		Soit $\overrightarrow{u} = \begin{pmatrix} {x} \\ {y} \end{pmatrix}$ un vecteur du plan : sa norme (notée $\Vert \overrightarrow{u} \Vert$) vaut $\Vert \overrightarrow{u} \Vert = \sqrt{\overrightarrow{u} \cdot \overrightarrow{u}} = \sqrt{x^2 + y^2}$.
	\end{formula}

	\begin{tip}[Caractéristiques d'un vecteur]
		On rappelle qu'un vecteur possède 3 caractéristiques :
		\begin{itemize}
			\item Une \textbf{norme} (sa longueur, par exemple si $\overrightarrow{u} = \overrightarrow{AB}$ alors $\Vert \overrightarrow{u} \Vert = AB$)
			\item Un \textbf{sens} (exemple : ``de $A$ vers $B$'' ou ``de haut en bas'')
			\item Une \textbf{direction} (la direction de la droite que porte le vecteur, horizontale ou verticale par exemple)
		\end{itemize}
	\end{tip}

	\subsection{Calcul}

	Il existe plusieurs méthodes pour calculer le produit scalaire en fonction de la situation dans laquelle on se trouve.

	\begin{formula}[Calcul avec un angle]
		Soient $\overrightarrow{u}$, $\overrightarrow{v}$ deux vecteurs du plan et $\theta$ l'angle orienté entre les deux. On a :
		\[ \overrightarrow{u} \cdot \overrightarrow{v} = \Vert \overrightarrow{u} \Vert \times \Vert \overrightarrow{v} \Vert \times \cos(\theta) \]
	\end{formula}

	\begin{formula}[Calcul avec un projeté orthogonal]
		Soient $A$, $B$ et $C$ trois points distincts du plan. On pose $P$ le projeté orthogonal de $C$ sur $(AB)$. Alors :
		\begin{itemize}
			\item Si $P \in [AB)$ alors $\overrightarrow{AB} \cdot \overrightarrow{AC} = AB \times AP$
			\item Si $P \notin [AB)$ alors $\overrightarrow{AB} \cdot \overrightarrow{AC} = - AB \times AP$
		\end{itemize}
	\end{formula}

	Si on ne possède que les normes de nos vecteurs, il est possible d'utiliser la formule de polarisation.

	\begin{formula}[Formule de polarisation]
		Soient $\overrightarrow{u}$ et $\overrightarrow{v}$ deux vecteurs du plan. Alors :
		\[ \overrightarrow{u} \cdot \overrightarrow{v} = \frac{1}{2} \left(\Vert \overrightarrow{u} + \overrightarrow{v} \Vert^2 - \Vert \overrightarrow{u} \Vert^2 - \Vert \overrightarrow{v} \Vert^2\right) \]
	\end{formula}

	\begin{tip}[Utilisation des formules]
		Il faut vraiment trouver la formule à utiliser selon l'énoncé de l'exercice.
		\newpar
		Par exemple, si on se trouve dans un repère et que l'on a les coordonnées des vecteurs, on pourra utiliser la formule de la \hyperref[definition]{définition}. À l'inverse, si on ne possède pas les coordonnées de nos vecteurs mais que l'on possède leur normes, il est possible d'utiliser la formule de polarisation.
		\newpar
		Voici un tableau récapitulatif pour $\overrightarrow{u}$ et $\overrightarrow{v}$ vecteurs du plan :
		\newpar
		\begin{whitetabularx}{|X|X|X|}
			\hline
			\textbf{Données} & \textbf{Formule} & \textbf{À utiliser si on possède...} \\
			\hline
			$\overrightarrow{u} = \begin{pmatrix} {x_1} \\ {y_1} \end{pmatrix}$ \medskip $\overrightarrow{v} = \begin{pmatrix} {x_2} \\ {y_2} \end{pmatrix}$. & $\overrightarrow{u} \cdot \overrightarrow{v} = x_1 \times x_2 + y_1 \times y_2$ \medskip (Calcul à partir des coordonnées.) & Les coordonnées de $\overrightarrow{u}$ et $\overrightarrow{v}$. \\
			\hline
			$\theta$ est l'angle orienté entre $\overrightarrow{u}$ et $\overrightarrow{v}$. & $\overrightarrow{u} \cdot \overrightarrow{v} = \Vert \overrightarrow{u} \Vert \times \Vert \overrightarrow{v} \Vert \times \cos(\theta)$ \medskip (Calcul à partir des normes et d'un angle.) & La norme de $\overrightarrow{u}$, la norme de $\overrightarrow{v}$ et l'angle $\theta$ entre les deux vecteurs. \\
			\hline
			$A$ et $B$ sont les deux extrémités de $\overrightarrow{u}$, $A$ et $C$ sont les deux extrémités de $\overrightarrow{v}$, et $P$ est le projeté orthogonal de $C$ sur $(AB)$. & $\overrightarrow{u} \cdot \overrightarrow{v} = \overrightarrow{AB} \cdot \overrightarrow{AC} = \pm AB \times AP$ \medskip $+$ si $P \in [AB)$ et $-$ sinon. \medskip (Calcul à partir d'une projection orthogonale.) & 3 points distincts (qui sont ici $A$, $B$ et $C$). \\
			\hline
			& $\overrightarrow{u} \cdot \overrightarrow{v} = \frac{\Vert \overrightarrow{u} + \overrightarrow{v} \Vert^2 - \Vert \overrightarrow{u} \Vert^2 - \Vert \overrightarrow{v} \Vert^2}{2}$ \medskip (Calcul à partir des normes.) & On possède la norme de $\overrightarrow{u}$, celle de $\overrightarrow{v}$ mais surtout celle de $\overrightarrow{u} + \overrightarrow{v}$. \\
			\hline
		\end{whitetabularx}
	\end{tip}

	\subsection{Théorème d'Al-Kashi}

	Le \textbf{théorème d'Al-Kashi} permet de calculer la longueur des côtés de n'importe quel triangle, qu'il soit rectangle ou non. Ainsi,

	\begin{formula}[Théorème d'Al-Kashi]
		Soient $A$, $B$ et $C$ trois points du plan non alignés (formant donc un triangle). On pose $a = BC$, $b = CA$ et $c = AB$. Alors :
		\[ c^2 = a^2 + b^2 - 2 \times a \times b \times \cos(\widehat{ACB}) \]
	\end{formula}

	\begin{demonstration}[Théorème d'Al-Kashi]
		En reprenant les notations de l'énoncé :
		\begin{align*}
			c^2 &= \Vert \overrightarrow{AB} \Vert^2 \\
			&= \Vert \overrightarrow{CB} - \overrightarrow{CA} \Vert^2 \text{ (par la relation de Chasles)} \\
			&= \Vert \overrightarrow{CB} \Vert^2 - 2(\overrightarrow{CB} \cdot \overrightarrow{CA}) + \Vert \overrightarrow{CA} \Vert^2 \text{ (par la formule de polarisation)} \\
			&= CB^2 - 2(CB \times CA \times \cos(\widehat{ACB})) + CA^2 \\
			&= a^2 + b^2 - 2 \times a \times b \times \cos(\widehat{ACB})
		\end{align*}
	\end{demonstration}

	\section{Géométrie}

	\subsection{Équation cartésienne d'une droite}

	\begin{formula}[Définition]
		Il est possible de décrire tous les points appartenant à une droite $\mathcal{D}$ par une équation appelée \textbf{équation cartésienne}.
		\newpar
		Une équation cartésienne de $\mathcal{D}$ est de la forme $ax + by + c = 0$ avec $a \neq 0$, $b \neq 0$ et $c$ réels, et où $x$ et $y$ sont des coordonnées de points.
	\end{formula}

	\begin{tip}
		Il est très facile de dire si oui ou non un point appartient à une droite si l'on possède l'équation cartésienne de cette droite.
		\newpar
		Par exemple, on définit la droite $\mathcal{D}$ par l'équation $y = x - 1$.
		\newpar
		Est-ce-que $A = (0; 1)$ appartient à $\mathcal{D}$ ? Remplaçons $x$ et $y$ par les coordonnées de $A$ : $1 = -1$ : c'est faux donc $A$ n'appartient pas à $\mathcal{D}$ car les coordonnées de $A$ ne vérifient par l'équation cartésienne de $\mathcal{D}$.
		\newpar
		Est-ce-que $B = (4; 3)$ appartient à $\mathcal{D}$ ? Remplaçons $x$ et $y$ par les coordonnées de $B$ : $3 = 3$ : c'est vrai donc $B$ appartient à $\mathcal{D}$ car les coordonnées de $B$ vérifient l'équation cartésienne de $\mathcal{D}$.
	\end{tip}

	\subsection{Vecteurs directeurs d'une droite}

	\begin{formula}[Définition]
		Soient $\mathcal{D}$ une droite et $\overrightarrow{u}$ un vecteur du plan non nul. Alors $\overrightarrow{u}$ est un \textbf{vecteur directeur} de $\mathcal{D}$ s'il existe deux points $A$ et $B$ appartenants à $\mathcal{D}$ et tels que $\overrightarrow{u} = \overrightarrow{AB}$.
	\end{formula}

	\includerepresentation{kccrxa9q}
	De plus, on a la propriété suivante qui peut s'avérer très utile :

	\begin{formula}[Colinéarité des vecteurs directeurs]
		$\overrightarrow{v}$ est un vecteur directeur de $\mathcal{D}$ si et seulement s'il est colinéaire au vecteur $\overrightarrow{u}$ précédent.
	\end{formula}

	Tous les vecteurs directeurs d'une droite sont donc colinéaires entre eux.

	\begin{tip}[Exemple]
		Soit $\mathcal{D}$ la droite définie par l'équation $y = 2x + 1$, montrons que $\overrightarrow{v} = \begin{pmatrix} 2 \\ 4\end{pmatrix}$ est un vecteur directeur de $\mathcal{D}$.
		\newpar
		Prenons deux points au hasard situés sur cette droite :
		\newline
		$x = 0$ donne $y = 1$, donc le point $A = (0; 1)$ appartient à $\mathcal{D}$.
		\newline
		$x = 1$ donne $y = 3$, donc le point $B = (1; 3)$ appartient à $\mathcal{D}$.
		\newpar
		Ainsi, un vecteur directeur de $\mathcal{D}$ est $\overrightarrow{u} = \overrightarrow{AB} = \begin{pmatrix} 1-0 \\ 3-1\end{pmatrix} = \begin{pmatrix} 1 \\ 2\end{pmatrix}$.
		\newpar
		Il reste à vérifier que $\overrightarrow{u}$ et $\overrightarrow{v}$ sont bien colinéaires, pour cela on peut utiliser la formule vue en seconde :
		\newline
		$2 \times 2 - 1 \times 4 = 0$ : $\overrightarrow{u}$ et $\overrightarrow{v}$ sont bien colinéaires et donc $\overrightarrow{v}$ est un vecteur directeur de $\mathcal{D}$.
	\end{tip}

	Il est facile de trouver un vecteur directeur d'une droite dont on connaît l'équation cartésienne.

	\begin{formula}[Coordonnées d'un vecteur directeur]
		Soit $\mathcal{D}$ une droite définie par l'équation $ax + by + c = 0$. Alors $\overrightarrow{u} = \begin{pmatrix} -b \\ a\end{pmatrix}$ est un vecteur directeur de $\mathcal{D}$.
	\end{formula}

	\begin{tip}[Exemple]
		Déterminons l'équation cartésienne de la droite $\mathcal{D}$ de vecteur directeur $\overrightarrow{u} = \begin{pmatrix} 1 \\ 2\end{pmatrix}$ et passant par $A = (1; 0)$.
		\newpar
		On a déjà $a$ et $b$ par la propriété précédente :
		\[ -b = 1 \iff b = -1 \text{ et } a = 2 \]
		Une équation cartésienne de la droite est $2x - y + c = 0$. Il reste à trouver $c$. Mais comme $\mathcal{D}$ passe par $A$, les coordonnées de $A$ vérifient l'équation cartésienne de $\mathcal{D}$.
		\newpar
		Remplaçons $x$ et $y$ par les coordonnées de $A$ dans l'équation cartésienne : \[ 2 + c = 0 \iff c = -2 \]
		L'équation cartésienne recherchée est donc $2x - y - 2 = 0$ ou encore $y = 2x - 2$.
	\end{tip}

	\begin{formula}[Propriétés]
		Soient $\mathcal{D}_1$ et $\mathcal{D}_2$ deux droites respectivement de vecteurs directeurs $\overrightarrow{u}$ et $\overrightarrow{v}$. Alors :
		\begin{itemize}
			\item $\mathcal{D}_1$ est parallèle à $\mathcal{D}_2$ si et seulement si $\overrightarrow{u}$ et $\overrightarrow{v}$ sont colinéaires.
			\item $\mathcal{D}_1$ est perpendiculaire à $\mathcal{D}_2$ si et seulement si $\overrightarrow{u} \cdot \overrightarrow{v} = 0$.
		\end{itemize}
	\end{formula}

	\begin{formula}[Orthogonalité]
		Si $\overrightarrow{u} \cdot \overrightarrow{v} = 0$ alors $\overrightarrow{u}$ et $\overrightarrow{v}$ sont dits \textbf{orthogonaux}.
	\end{formula}

	\subsection{Vecteurs normaux à une droite}

	\begin{formula}[Définition]
		Soient $\mathcal{D}$ une droite de vecteur directeur $\overrightarrow{u}$ et $\overrightarrow{n}$ un vecteur du plan non nul. Alors $\overrightarrow{n}$ est un \textbf{vecteur normal} à $\mathcal{D}$ si $\overrightarrow{u}$ et $\overrightarrow{n}$ sont orthogonaux entre-eux.
	\end{formula}

	\includerepresentation{ydshrp8y}

	De plus, on a la propriété suivante qui peut s'avérer très utile :

	\begin{formula}[Colinéarité des vecteurs normaux]
		$\overrightarrow{m}$ est un vecteur normal à $\mathcal{D}$ si et seulement s'il est colinéaire au vecteur $\overrightarrow{n}$ précédent.
	\end{formula}

	Tous les vecteurs normaux d'une droite sont donc colinéaires entre-eux.

	Il est facile de trouver un vecteur normal à une droite dont on connaît l'équation cartésienne.

	\begin{formula}[Coordonnées d'un vecteur normal]
		Soit $\mathcal{D}$ une droite définie par l'équation $ax + by + c = 0$. Alors $\overrightarrow{n} = \begin{pmatrix} a \\ b\end{pmatrix}$ est un vecteur normal à $\mathcal{D}$.
	\end{formula}

	Soient $\mathcal{D}_1$ et $\mathcal{D}_2$ deux droites respectivement de vecteurs directeurs $\overrightarrow{u}$ et $\overrightarrow{v}$. Alors :

	\begin{formula}
		$\mathcal{D}_1$ est perpendiculaire à $\mathcal{D}_2$ si et seulement si $\overrightarrow{u}$ est normal à $\mathcal{D}_2$.
	\end{formula}

	\begin{tip}[Exemple]
		Déterminons l'équation cartésienne de la droite $\mathcal{D}$ admettant pour vecteur normal $\overrightarrow{n} = \begin{pmatrix} -1 \\ -1\end{pmatrix}$ et passant par l'origine $O = (0; 0)$.
		\newpar
		On a déjà $a$ et $b$ par la propriété précédente :
		\newline
		$a = -1$
		\newline
		$b = -1$
		\newpar
		Une équation cartésienne de la droite est $-x - y + c = 0$. Il reste à trouver $c$. Mais comme $\mathcal{D}$ passe par l'origine, les coordonnées de $O$ vérifient l'équation cartésienne de $\mathcal{D}$.
		\newpar
		Remplaçons $x$ et $y$ par les coordonnées de $O$ dans l'équation cartésienne : $c = 0$.
		L'équation cartésienne recherchée est donc $-x - y = 0$ ou encore $y = -x$.
	\end{tip}

	\subsection{Description d'un cercle}

	De la même manière que pour les droites, il est possible de décrire l'ensemble des points appartenant à un cercle à l'aide d'une équation.

	\begin{formula}[Description par équation cartésienne]
		Soit $\mathcal{C}$ un cercle de centre $O = (x_O; y_O)$ et de rayon $R$.
		\newpar
		Une équation cartésienne de $\mathcal{C}$ est de la forme $(x - x_O)^2 + (y - y_O)^2 = R^2$ avec $x$ et $y$ qui sont des coordonnées de points.
	\end{formula}

	On peut de manière équivalente, décrire un cercle à l'aide du produit scalaire.

	\begin{formula}[Description par produit scalaire]
		Soient $A$ et $B$ deux points du plan. Alors l'ensemble des points $M$ tels que $\overrightarrow{MA} \cdot \overrightarrow{MB} = 0$ est le cercle de diamètre $[AB]$.
	\end{formula}

	\begin{demonstration}[Description par produit scalaire]
		On pose $A = (x_A; y_A)$, $B = (x_B; y_B)$ et on cherche les points $M = (x; y)$ tels que $\overrightarrow{MA} \cdot \overrightarrow{MB} = 0$.
		\newpar
		Soit $O$ le milieu de $[AB]$ :
		\newpar
		$\overrightarrow{MA} \cdot \overrightarrow{MB} = 0$
		\newline
		$\iff (\overrightarrow{MO} + \overrightarrow{OA}) \cdot (\overrightarrow{MO} + \overrightarrow{OB}) = 0$
		\newline
		$\iff (\overrightarrow{MO} + \overrightarrow{OA}) \cdot (\overrightarrow{MO} - \overrightarrow{OA}) = 0$
		\newline
		$\iff (\overrightarrow{MO} \cdot \overrightarrow{MO}) - (\overrightarrow{OA} \cdot \overrightarrow{OA}) = 0$
		\newline
		$\iff MO^2 - OA^2 = 0$
		\newline
		$\iff MO = OA$
		\newpar
		Donc l'ensemble cherché est l'ensemble des points situés à une distance $OA$ du point $O$, c'est bien le cercle de centre $O$ et de diamètre $[AB]$.
	\end{demonstration}

	En réalité, les deux points précédents sont deux manières différentes de décrire un cercle.
	%</content>
\end{document}
